\section*{Bab \ref{ch:g12:condprob} --- Peluang Bersyarat dan Kebebasan}

\begin{solution}{exo:g12:condprob:1}
Ada $12$ kartu bergambar, dan $4$ di antaranya raja:
$\pcond{\text{bergambar}}{\text{raja}} = \frac{4}{12} = \frac13$.

\omterm{def:g12:condprob:indep}{Kebebasannya}: $\P(\text{raja} \cap \text{hati}) = \frac{1}{52}$ (yaitu raja
hati), sedangkan
$\P(\text{raja})\,\P(\text{hati}) = \frac{4}{52}\times\frac{13}{52}
= \frac{1}{52}$. Keduanya sama: jadi \omterm{def:g11:prob:model}{kejadiannya} \omterm{def:g12:condprob:indep}{saling bebas}.
\end{solution}

\begin{solution}{exo:g12:condprob:2}
\emph{1.} Cabang pertamanya: merah $\frac58$, biru $\frac38$; cabang keduanya
(tanpa pengembalian): sesudah merah, merah $\frac47$ / biru $\frac37$; sesudah
biru, merah $\frac57$ / biru $\frac27$.

\emph{2.} $\P(RR) = \frac58 \times \frac47 = \frac{5}{14}$. Menurut peluang
totalnya,
\[
\P(\text{merah ke-2}) = \frac58\cdot\frac47 + \frac38\cdot\frac57
= \frac{20 + 15}{56} = \frac58 .
\]
(Sama dengan pengambilan pertamanya --- menurut kesetangkupannya, bola kedua adalah
bola yang acak seragam dari gucinya.)
\end{solution}

\begin{solution}{exo:g12:condprob:3}
$\P(A) = \frac{6}{36} = \frac16$ (ada enam pasangan yang berjumlah $7$),
$\P(B) = \frac16$, dan $A \cap B = \{(3,4)\}$ berpeluang
$\frac{1}{36} = \P(A)\P(B)$: jadi $A$ dan $B$ \omterm{def:g12:condprob:indep}{saling bebas}. (Jumlah $7$ itu
istimewa: apa pun dadu pertamanya, tepat satu nilai dadu kedua memberikannya.)

$\P(C) = \frac{5}{36}$ dan $B \cap C = \{(3,3)\}$, sehingga
$\P(B \cap C) = \frac{1}{36} \neq \frac16 \times \frac{5}{36}
= \frac{5}{216}$. Jadi tidak \omterm{def:g12:condprob:indep}{saling bebas}.
\end{solution}

\begin{solution}{exo:g12:condprob:4}
\emph{1.} Peluang total dengan partisi menurut mesinnya:
\[
\P(D) = 0.5\times0.01 + 0.3\times0.02 + 0.2\times0.04
= 0.005 + 0.006 + 0.008 = 0.019 = 1.9\% .
\]

\emph{2.} Bayes memberi
$\pcond{D}{M_3} = \dfrac{0.008}{0.019} = \dfrac{8}{19} \approx 0.42$. Jadi
mesin yang menghasilkan hanya seperlima keluarannya menyumbang lebih dari $40\%$
cacatnya.
\end{solution}

\begin{solution}{exo:g12:condprob:5}
Syaratnya adalah
\[
\frac{0.99\,p}{0.99\,p + 0.05\,(1-p)} \geq 0.9 .
\]
Penyebutnya positif, sehingga ini terbaca
$0.99p \geq 0.891p + 0.045 - 0.045p$, \emph{yaitu} $0.144\,p \geq 0.045$,
\emph{yaitu}
\[
p \geq \frac{0.045}{0.144} = 0.3125 .
\]
Ujinya sendirian mencapai kepastian $90\%$ hanya jika penyakitnya sudah menyerang
lebih dari $31\%$ penduduk yang diuji --- dan itulah sebabnya penapisan massal
bagi penyakit langka menuntut uji penegasan.
\end{solution}

\begin{solution}{exo:g12:condprob:6}
\emph{1.} Ada tiga \omterm{def:g12:seq:sequence}{barisan} yang mewujudkan $E$ (GGA, GAG, AGG), masing-masing berpeluang
$p^2(1-p)$ menurut \omterm{def:g12:condprob:indep}{kebebasannya}, jadi $\P(E) = 3p^2(1 - p)$.

\emph{2.} $f(p) = 3p^2 - 3p^3$ mempunyai $f'(p) = 6p - 9p^2 = 3p(2 - 3p)$,
yang positif lalu negatif pada $\intoo{0}{1}$: jadi \omterm{def:g10:functions:extrema}{maksimumnya} di $p = \frac23$, dan di sana
$\P(E) = 3\cdot\frac49\cdot\frac13 = \frac49$.
\end{solution}

\begin{solution}{exo:g12:condprob:7}
Katakanlah kamu memilih pintu 1 dan pembawa acaranya membuka pintu 3 (\omterm{def:g11:prob:model}{kejadian} $H_3$). Dengan
$B_i$ = ``hadiahnya di balik pintu $i$'', berlaku $\P(B_i) = \frac13$ dan
\[
\pcond{B_1}{H_3} = \frac12, \qquad
\pcond{B_2}{H_3} = 1, \qquad
\pcond{B_3}{H_3} = 0
\]
(jika hadiahnya di balik pintumu, pembawa acaranya memilih antara pintu 2 dan 3
secara acak; jika hadiahnya di balik pintu 2, ia terpaksa membuka pintu 3). Bayes:
\[
\pcond{H_3}{B_1}
= \frac{\frac13\cdot\frac12}{\frac13\cdot\frac12 + \frac13\cdot1 + 0}
= \frac{1/6}{1/2} = \frac13,
\qquad
\pcond{H_3}{B_2} = \frac{\frac13}{\frac12} = \frac23 .
\]
Bertahan menang dengan peluang $\frac13$, sedangkan berpindah dengan peluang
$\frac23$.
\end{solution}

\begin{solution}{exo:g12:condprob:8}
\emph{1.} Suara terbanyaknya keliru ketika sekurang-kurangnya dua dari tiga salinannya
terbalik:
\[
3\varepsilon^2(1-\varepsilon) + \varepsilon^3
= 3(0.01)(0.9) + 0.001 = 0.028,
\]
jauh lebih kecil daripada $\varepsilon = 0.1$: jadi penyandian pengulangan itu berhasil.

\emph{2.} Jika $1$ yang dikirim (sebagai $111$), menerima $101$ menuntut tepat satu
pembalikan: peluangnya $\varepsilon(1-\varepsilon)^2 = 0.081$. Jika $0$ yang dikirim
(sebagai $000$), diperlukan dua pembalikan: $\varepsilon^2(1-\varepsilon) = 0.009$.
Bayes dengan peluang awal yang sama:
\[
\pcond{101}{\,\text{terkirim }1} = \frac{0.081}{0.081 + 0.009} = 0.9 .
\]
Penerjemahan menurut suara terbanyak ($101 \mapsto 1$) memang terkaan yang lebih mungkin.
\end{solution}

\begin{solution}{pb:g12:condprob:1}
\textbf{1.} $\P(\text{merah}) = \frac12 \cdot \frac23 + \frac12
\cdot \frac14 = \frac13 + \frac18 = \frac{11}{24}$.

\textbf{2.} $\P(\text{I} \mid \text{merah}) =
\dfrac{\frac12 \cdot \frac23}{\frac{11}{24}} =
\dfrac{1/3}{11/24} = \dfrac{8}{11}$.

\textbf{3.} $\P(\text{genap} \cap \leq 4) = \P(\{2, 4\}) =
\frac13$ dan $\P(\text{genap})\P(\leq 4) = \frac12 \cdot
\frac23 = \frac13$: jadi \omterm{def:g12:condprob:indep}{saling bebas} (di luar dugaan!). Untuk
$\leq 3$: $\P(\{2\}) = \frac16$ terhadap
$\frac12 \cdot \frac12 = \frac14$: jadi bergantungan.

\textbf{4.} \omterm{def:g10:proba:operations}{Komplemen} \omterm{def:g11:prob:model}{kejadian}
yang \omterm{def:g12:condprob:indep}{saling
bebas} juga
\omterm{def:g12:condprob:indep}{saling bebas}, sehingga
\[
\P(\text{tak satu pun}) = 0.1 \times 0.1 = 0.01,
\qquad
\P(\text{sekurang-kurangnya satu}) = 0.99 .
\]

\textbf{5.} $\frac{4}{52} \times \frac{3}{51} =
\frac{1}{221}$.

\textbf{6.} Yang sakit: $100$; yang sehat: $99\,900$. Positifnya:
$99$ yang benar (99\,\% dari $100$) dan $4\,995$ yang palsu (5\,\% dari
$99\,900$), jadi $5\,094$ positif seluruhnya.

\textbf{7.} $\P(\text{sakit} \mid +) = \frac{99}{5\,094}
\approx 1.9\,\%$: jadi dari orang yang ditandai ujinya, kurang dari
satu dari lima puluh yang benar-benar sakit.

\textbf{8.} $\dfrac{0.001 \times 0.99}{0.001 \times 0.99 +
0.999 \times 0.05} = \dfrac{0.00099}{0.05094} \approx 0.019$:
jadi cacahnya dan rumusnya sepakat.

\textbf{9.} Yang diacaukan: $\P(+ \mid \text{sakit}) = 99\,\%$
(yaitu mutu ujinya yang diiklankan) dengan
$\P(\text{sakit} \mid +)$ (yaitu pertanyaan pasiennya yang sebenarnya). Dengan
laju dasar $\frac{1}{1000}$, orang sakit itu begitu jarang sehingga
bahkan bocoran $5\,\%$ dari mayoritas sehat yang sangat besar itu menenggelamkan
positif benarnya lima puluh berbanding satu.

\textbf{10.} $\dfrac{0.019 \times 0.99}{0.019 \times 0.99 +
0.981 \times 0.05} \approx 28\,\%$. Tiap positif yang \omterm{def:g12:condprob:indep}{saling bebas}
mengalikan peluangnya dengan faktor yang sama: buktinya bertimbun,
dan positif ketiga akan mendorongnya melewati $90\,\%$ --- dan itulah
sebabnya uji penegasan ada.

\textbf{11.} $\dfrac{0.05 \times 0.99}{0.05 \times 0.99 +
0.95 \times 0.05} \approx 51\,\%$: jadi pada kelompok berisiko tinggi
uji yang sama itu informatif. Pelajarannya: menapislah di tempat laju dasarnya
cukup besar; menapis massal penyakit yang langka justru memproduksi
alarm palsu sampai ribuan.

\textbf{12.} Harapan banyaknya kecocokan tak bersalah:
$10^7 \times 10^{-6} = 10$. Di antara kira-kira $11$ orang yang cocok
($10$ tak bersalah, $1$ bersalah --- jika pelakunya ada di
kota itu), tersangka yang ditemukan \emph{hanya lewat kecocokannya}
tak bersalah dengan peluang kira-kira $\frac{10}{11} \approx
91\,\%$.

\textbf{13.} Sang jaksa mengutip
$\P(\text{cocok} \mid \text{tak bersalah}) = 10^{-6}$; padahal pengadilan
memerlukan $\P(\text{tak bersalah} \mid \text{cocok})$ --- yaitu kekacauan pertanyaan
9 yang diucapkan di bawah sumpah. (Dan pada perkara nyata yang
memasyhurkan galat ini, kekeliruan kedua memperparahnya:
lihat pertanyaan 14.)

\textbf{14.} Pengkuadratan yang naif: $\left(\frac{1}{8500}\right)^2
\approx \frac{1}{72 \text{ juta}}$. Dengan musibah kedua
yang lima kali lebih mungkin jika yang pertama diketahui:
$\frac{1}{8500} \times \frac{5}{8500} \approx
\frac{1}{14.5 \text{ juta}}$ --- lima kali lebih besar, dan
itu pun masih mengabaikan bahwa hipotesis \emph{tandingannya} harus
ditimbang oleh Bayes juga. Pelajarannya: mengalikan peluang
menuntut \omterm{def:g12:condprob:indep}{kebebasan}, dan sebab yang dikongsi (genetika,
lingkungan) merusaknya.

\textbf{15.} Bayes mengalikan nisbah kebolehjadian dari
\emph{semua} buktinya menjadi satu peluang akhir: kecocokannya, alibinya,
seratnya, masing-masing menggeser peluangnya. Kalimat bagi anggota juri:
\emph{tanyakanlah seberapa mungkin buktinya di bawah kedua
hipotesisnya, dan jangan sekali-kali menukar satu \omterm{def:g12:condprob:cond}{peluang bersyarat} dengan
kebalikannya}.

\textbf{16.} $\dfrac{0.4 \times 0.15}{0.4 \times 0.15 +
0.6 \times 0.005} = \dfrac{0.060}{0.063} \approx 95\,\%$
berupa sampah: jadi satu kata saja nyaris memutuskannya.

\textbf{17.} Lima salinan gagal ketika $3$, $4$ atau $5$ terbalik:
$\binom53 \varepsilon^3(1-\varepsilon)^2 + \binom54
\varepsilon^4 (1 - \varepsilon) + \varepsilon^5 \approx
0.0086$ --- terhadap $0.028$ untuk tiga salinan: jadi pengulangan
membeli keandalan dengan harga lebar pitanya.

\textbf{18.} Pembalikan yang genap: $\binom40 0.9^4 + \binom42 0.9^2
0.1^2 + \binom44 0.1^4 = 0.6561 + 0.0486 + 0.0001 = 0.7048$;
dan $\frac{1 + 0.8^4}{2} = \frac{1 + 0.4096}{2} = 0.7048$:
jadi rumus muslihatnya melacak biasnya $(1 - 2\varepsilon)$
lewat tiap penerusnya.

\textbf{19.} Kata berbau sampah berjalan bergerombol --- pesan dengan
``undian'' besar kemungkinan memuat ``menang'', sehingga bukti gabungannya
lebih lemah daripada yang dipura-purakan hasil kalinya. Penyaringnya
tetap bekerja sebab pembesar-besarannya biasanya mendorong ke
\emph{arah yang benar}: penggolongan hanya perlu sisi yang benar
dari $50\,\%$, bukan peluang akhir yang benar.

\textbf{20.} Klinik: penyakit yang langka mengubah uji yang kuat
menjadi vonis yang lemah --- perbaruilah dari laju dasarnya, lalu
timbunlah. Ruang sidang: kecocokan langka di kota besar mendakwa
kotanya, bukan orangnya --- kebolehjadian bukanlah vonis.
Kotak masuk: banyak petunjuk lemah, yang dikalikan, memilah suratnya --- bahkan
secara naif. Ketiga jebakannya: lupakan peluang awalmu dan kamu tertipu;
palsukan \omterm{def:g12:condprob:indep}{kebebasannya} dan kamu tertipu berat; bacalah bukti sepenggal-sepenggal
dan kamu tak pernah tahu apa yang kamu ketahui.

\end{solution}
